\documentclass{notaaula}

        \titulonota{Som e efeito Doppler}
        \creditos{Turma 2A · Física · Dezembro/2026 · Prof. Josué Cavalcante}

        \begin{document}
        \cabecalho

        \begin{sobre}
        Acústica: produção e propagação do som, qualidades sonoras, intensidade e efeito Doppler. Habilidades em foco: EM13CNT106, EM13CNT107 e EM13CNT310.
        \end{sobre}

        \section{Produção e propagação do som}

\begin{copiar}{Onda sonora}
O som é uma onda mecânica produzida por vibração. No ar, propaga-se como compressões e rarefações
longitudinais; não se propaga no vácuo. Vale a relação
\[
  v=\lambda f.
\]
No ar a $20^\circ\mathrm C$, usamos aproximadamente $v=340\un{m/s}$. Em geral, a rapidez é maior
em sólidos do que em líquidos e gases, pois depende das propriedades elásticas e da densidade do meio.
\end{copiar}

\begin{exemplo}
Um som de $680\un{Hz}$ propaga-se no ar a $340\un{m/s}$. Seu comprimento de onda é
\[
  \lambda=\frac{v}{f}=\frac{340}{680}
  =\resultado{\dec{0,50}\un{m}}.
\]
\end{exemplo}

\section{Qualidades sonoras}

\begin{copiar}{Altura, intensidade e timbre}
\begin{itens}
\item \dest{Altura:} sensação de grave ou agudo; relaciona-se principalmente à frequência.
\item \dest{Intensidade:} energia por área e por tempo; relaciona-se à amplitude e à sensação de volume.
\item \dest{Timbre:} distingue fontes que emitem a mesma nota; depende da forma da onda e dos harmônicos.
\end{itens}
A faixa audível humana típica vai de cerca de $20\un{Hz}$ a $20\,000\un{Hz}$, variando com a pessoa e a idade.
\atencao{Som alto no cotidiano pode significar intenso; em Física, altura sonora significa frequência.}
\end{copiar}

\begin{copiar}{Nível sonoro}
A intensidade é $I=P/A$. O nível sonoro em decibéis é
\[
  \beta=10\log\left(\frac{I}{I_0}\right),
  \qquad I_0=10^{-12}\un{W/m^2}.
\]
A escala é logarítmica: multiplicar a intensidade por 10 acrescenta $10\un{dB}$.
Exposição prolongada a níveis elevados pode causar perda auditiva.
\end{copiar}

\section{Eco, reverberação e ressonância}

\begin{copiar}{Acústica de ambientes}
Eco é uma reflexão percebida separadamente do som original. Reverberação é a persistência causada
por muitas reflexões próximas. Materiais porosos absorvem parte da energia sonora; superfícies duras
refletem mais. Instrumentos musicais usam ressonância para reforçar frequências naturais.
\end{copiar}

\section{Efeito Doppler}

\begin{copiar}{Movimento relativo}
Quando fonte e observador se aproximam, a frequência percebida aumenta; quando se afastam, diminui.
Para velocidades ao longo da linha que une fonte e observador,
\[
  f'=f\frac{v\pm v_o}{v\mp v_s}.
\]
Escolha os sinais para que aproximação aumente $f'$ e afastamento diminua $f'$.
\simbolos{$v_o$ é a rapidez do observador; $v_s$, a da fonte; $v$, a do som.}
\end{copiar}

\begin{exemplo}[Fonte aproximando-se]
Uma sirene de $600\un{Hz}$ aproxima-se de observador parado a $34\un{m/s}$, com $v=340\un{m/s}$.
\[
  f'=600\frac{340}{340-34}
  \approx\resultado{667\un{Hz}}.
\]
\end{exemplo}

\begin{dica}
O efeito Doppler muda a frequência percebida, não a frequência de emissão da fonte. Para observador
parado, cristas ficam comprimidas à frente da fonte e espaçadas atrás dela.
\end{dica}

\begin{exercicios}
\nivel{1}{Conceitos}
\begin{questoes}
\item Por que o som não se propaga no vácuo?
\item Relacione frequência às sensações de grave e agudo.
\item O que diferencia o timbre de duas fontes?
\item Diferencie eco de reverberação.
\item Na aproximação entre fonte e observador, $f'$ aumenta ou diminui?
\nivel{2}{Aplicação}
\item Um som de $170\un{Hz}$ propaga-se a $340\un{m/s}$. Ache $\lambda$.
\item Uma onda sonora tem $\lambda=\dec{0,85}\un{m}$ no ar. Ache $f$.
\item Multiplicar a intensidade por 100 aumenta o nível em quantos decibéis?
\item Um eco retorna $\dec{0,40}\un{s}$ após a emissão. Com $v=340\un{m/s}$, ache a distância da parede.
\nivel{3}{Síntese}
\item Uma sirene de $500\un{Hz}$ aproxima-se a $20\un{m/s}$ de observador parado. Use $v=340\un{m/s}$ e ache $f'$.
\item Explique por que ambulâncias parecem mais agudas quando se aproximam.
\item Proponha duas medidas para reduzir reverberação em uma sala.
\end{questoes}
\gabarito{5) aumenta; 6) $2\un{m}$; 7) $400\un{Hz}$; 8) $20\un{dB}$;
9) $68\un{m}$; 10) aproximadamente $531\un{Hz}$.}
\end{exercicios}


\end{document}
